\documentclass{article}
\usepackage[T1]{fontenc}
\usepackage{lmodern}
\usepackage{graphicx}
\usepackage{amsmath,amssymb}
\usepackage[a4paper,margin=2cm]{geometry}
\usepackage[absolute,overlay]{textpos}
\setlength{\TPHorizModule}{1mm}
\setlength{\TPVertModule}{1mm}
\usepackage[indonesian]{babel}
\usepackage{hyperref}
\setlength{\emergencystretch}{2em}
% unit: Halaman 87

\begin{document}

% === Halaman 87 (llm) ===

\begin{itemize}
  \item Contoh:
  \end{itemize}

Plainteks: \texttt{kripto} $(10 \ 17 \ 8 \ 15 \ 19 \ 14)$

$n = 26$, ambil $m = 7$ (7 relatif prima dengan 26)

\vspace{5mm}

Enkripsi: $C \equiv 7P + 10 \pmod{26}$

\begin{align*}
 p_1 = 10 & \to c_1 \equiv 7 \cdot 10 + 10 \equiv 80 \equiv 2 \pmod{26} && (\text{huruf 'C'})
 p_2 = 17 & \to c_2 \equiv 7 \cdot 17 + 10 \equiv 129 \equiv 25 \pmod{26} && (\text{huruf 'Z'})
 p_3 = 8 & \to c_3 \equiv 7 \cdot 8 + 10 \equiv 66 \equiv 14 \pmod{26} && (\text{huruf 'O'})
 p_4 = 15 & \to c_4 \equiv 7 \cdot 15 + 10 \equiv 115 \equiv 11 \pmod{26} && (\text{huruf 'L'})
 p_5 = 19 & \to c_5 \equiv 7 \cdot 19 + 10 \equiv 143 \equiv 13 \pmod{26} && (\text{huruf 'N'})
 p_6 = 14 & \to c_6 \equiv 7 \cdot 14 + 10 \equiv 108 \equiv 4 \pmod{26} && (\text{huruf 'E'})
\end{align*}

\vspace{5mm}

Cipherteks: \texttt{CZOLNE}

\end{document}
